%% ==========================================================================
%%  Approach 13 --- Conjecture 2 for an arbitrary number of agents.
%%
%%  Self-contained. External input: Theorem 5.1 of [TWYZ23] together with the
%%  terminal behaviour of its Algorithm 3. Everything after the terminal
%%  partial allocation is proved here. No enumeration is used anywhere.
%%
%%  REFERENCING: theorem-like environments share one counter in this report,
%%  so \Cref renders all of them as "Theorem". Use explicit
%%  "Lemma~\ref{...}", "Theorem~\ref{...}", "Corollary~\ref{...}".
%% ==========================================================================
\section{Approach 13: Unit Subsidies for an Arbitrary Number of Agents}
\label{sec:approach13}

This section proves Conjecture~\ref{conj:main} for every number of agents.
All notation used is defined below; the only external input is
Theorem~\ref{thm:g3-twyz} and the description of its algorithm in
\Cref{sec:g3-terminal}.

\subsection{Setting}
\label{sec:g3-setting}

Let $\agents = \set{1,\dots,n}$ be a set of agents and $\items$ a finite set
of indivisible chores. Each agent $i \in \agents$ has a \emph{dichotomous}
cost function $\cost_i : 2^{\items} \to \Z_{\ge 0}$, meaning
\begin{equation}
\label{eq:g3-dich}
  \cost_i(\emptyset) = 0,
  \qquad
  \cost_i(S \cup \set e) - \cost_i(S) \in \set{0,1}
  \quad\text{for all } S \subseteq \items,\ e \in \items \setminus S .
\end{equation}
Write $\cost_i(e \mid S) := \cost_i(S \cup \set e) - \cost_i(S)$.

\begin{observation}
\label{obs:g3-basic}
Every dichotomous $\cost_i$ is integer valued and non-decreasing: for
$S \subseteq T \subseteq \items$, $\cost_i(S) \le \cost_i(T)$. Consequently,
for any $S,T$,
\begin{equation}
\label{eq:g3-gap}
  \cost_i(S) < \cost_i(T) \quad\Longrightarrow\quad \cost_i(S) \le \cost_i(T) - 1 .
\end{equation}
\end{observation}
\begin{proof}
Enumerate $T \setminus S = \set{f_1,\dots,f_q}$ and telescope: $\cost_i(T) -
\cost_i(S) = \sum_{u=1}^{q} \cost_i(f_u \mid S \cup \set{f_1,\dots,f_{u-1}})$,
a sum of terms in $\set{0,1}$ by \eqref{eq:g3-dich}. Taking $S = \emptyset$
and using $\cost_i(\emptyset)=0$ shows every value is a non-negative integer;
the displayed sum is non-negative, giving monotonicity. Statement
\eqref{eq:g3-gap} is then the integrality of $\cost_i(T) - \cost_i(S)$.
\end{proof}

A \emph{partial allocation} is a tuple $\alloc = (\bundle1,\dots,\bundle n)$
of pairwise disjoint subsets of $\items$; its \emph{residue} is
$R(\alloc) := \items \setminus \bigcup_{i} \bundle i$. It is an
\emph{allocation} when $R(\alloc) = \emptyset$. A partial allocation is
\emph{envy-free} if
\begin{equation}
\label{eq:g3-pef}
  \cost_i(\bundle i) \le \cost_i(\bundle j) \qquad \text{for all } i,j \in \agents .
\end{equation}
For $\subsidy \in \Z_{\ge0}^n$, the pair $(\alloc,\subsidy)$ is an
\emph{envy-free solution} if
\begin{equation}
\label{eq:g3-ef}
  \cost_i(\bundle i) - \subsidy_i \le \cost_i(\bundle j) - \subsidy_j
  \qquad \text{for all } i,j \in \agents .
\end{equation}
Conjecture~\ref{conj:main} asserts that for every instance there are an
allocation $\alloc$ and a subsidy $\subsidy \in \set{0,1}^n$, with
$\sum_i \subsidy_i \le n-1$, such that \eqref{eq:g3-ef} holds, computable in
polynomial time.

\subsection{The terminal state of the external algorithm}
\label{sec:g3-terminal}

\begin{theorem}[\cite{TWYZ23}, Theorem 5.1]
\label{thm:g3-twyz}
For every $n$ and all dichotomous $\cost_1,\dots,\cost_n$ there is a
polynomial-time algorithm returning an envy-free partial allocation $\alloc$
with $\abs{R(\alloc)} \le n-1$.
\end{theorem}

The algorithm of Theorem~\ref{thm:g3-twyz} maintains an envy-free partial
allocation $X = (X_1,\dots,X_n)$ together with the \emph{equality graph}
\begin{equation}
\label{eq:g3-graph}
  G = (\agents, E),
  \qquad
  (i,j) \in E \iff i \ne j \ \text{ and } \ \cost_i(X_i) = \cost_i(X_j),
\end{equation}
and repeats the following three rules in the stated order, halting only when
no rule applies.
\begin{enumerate}
  \item[(R1)] If some $e \in R(X)$ and some $i \in \agents$ satisfy
        $\cost_i(e \mid X_i) = 0$, assign $e$ to $i$.
  \item[(R2)] If some $e \in R(X)$ and some arc $(i,j) \in E$ lying on a
        directed cycle of $G$ satisfy $\cost_i(e \mid X_j) = 0$, rotate the
        bundles along that cycle and assign $e$ to $i$.
  \item[(R3)] Otherwise select a strongly connected component $S$ of $G$ with
        no arc of $E$ leaving $S$. If $\abs{R(X)} \ge \abs S$, assign one
        residual chore to each agent of $S$; otherwise halt.
\end{enumerate}
The loop runs while $R(X) \ne \emptyset$, and the only exit with
$R(X) \ne \emptyset$ is the final clause of (R3).

For the remainder of this section fix the following data. Let $X$ be the
partial allocation at which the algorithm halts, let $R := R(X)$ and
$r := \abs R$, and, if $r \ge 1$, let $S$ be the component selected by (R3)
at the halting step. All of $\alloc$'s envy-freeness, $G$, $R$, $S$ refer to
this terminal state.

\begin{lemma}
\label{lem:g3-terminal}
Suppose $r \ge 1$. Then
\begin{enumerate}
  \item[(a)] $r < \abs S$;
  \item[(b)] $\cost_i(e \mid X_i) = 1$ for every $i \in \agents$ and every
        $e \in R$;
  \item[(c)] $\cost_i(e \mid X_j) = 1$ for every arc $(i,j) \in E$ with
        $i,j \in S$, and every $e \in R$;
  \item[(d)] $\cost_i(X_i) < \cost_i(X_j)$ for every $i \in S$ and every
        $j \in \agents \setminus S$.
\end{enumerate}
\end{lemma}
\begin{proof}
Since $r \ge 1$ the algorithm halted at the final clause of (R3), so neither
(R1) nor (R2) applied at the halting state, and $\abs{R} < \abs S$, which is
(a).

(b) If $\cost_i(e \mid X_i) = 0$ for some $i$ and some $e \in R$, rule (R1)
would apply. By \eqref{eq:g3-dich} the marginal lies in $\set{0,1}$, hence
equals $1$.

(c) Let $(i,j) \in E$ with $i,j \in S$. As $S$ is strongly connected there is
a directed path in $G$ from $j$ to $i$; choosing one with no repeated vertex
and appending the arc $(i,j)$ produces a directed cycle of $G$ containing
$(i,j)$. Were $\cost_i(e \mid X_j) = 0$ for some $e \in R$, rule (R2) would
apply to that arc and cycle. Hence the marginal is not $0$, so it is $1$.

(d) Let $i \in S$ and $j \notin S$. No arc of $E$ leaves $S$, so
$(i,j) \notin E$, i.e.\ $\cost_i(X_i) \ne \cost_i(X_j)$. Envy-freeness
\eqref{eq:g3-pef} gives $\cost_i(X_i) \le \cost_i(X_j)$, whence the
inequality is strict.
\end{proof}

\subsection{The completion}
\label{sec:g3-construction}

Assume $r \ge 1$ and write $R = \set{e_1,\dots,e_r}$. By
Lemma~\ref{lem:g3-terminal}(a) we have $r < \abs S$, so we may fix $r$
distinct agents
\[
  T = \set{t_1,\dots,t_r} \subseteq S ,
\]
and we do so arbitrarily. Define the allocation $\alloc = (A_1,\dots,A_n)$ by
\begin{equation}
\label{eq:g3-alloc}
  A_i :=
  \begin{cases}
    X_{t_k} \cup \set{e_k}, & i = t_k \ \ (1 \le k \le r),\\[2pt]
    X_i, & i \in \agents \setminus T .
  \end{cases}
\end{equation}
Since the $t_k$ are distinct and every chore of $R$ is used exactly once,
$\alloc$ is an allocation of $\items$. Members of $T$ are called
\emph{recipients}.

\begin{definition}[Subsidy set]
\label{def:g3-P}
Let $P \subseteq \agents$ be the smallest set satisfying
\begin{enumerate}
  \item[(P1)] $T \subseteq P$; and
  \item[(P2)] if $i \in \agents \setminus S$ and $(i,j) \in E$ for some
        $j \in P$, then $i \in P$.
\end{enumerate}
Equivalently, $P$ is obtained by initialising $P := T$ and repeatedly
adjoining any $i \notin S$ having an arc of $E$ into the current set, until
no further agent can be adjoined; the process terminates because $\agents$ is
finite and $P$ only grows. Define
\begin{equation}
\label{eq:g3-subsidy}
  \subsidy_i := \begin{cases} 1, & i \in P,\\ 0, & i \notin P.\end{cases}
\end{equation}
\end{definition}

\begin{lemma}
\label{lem:g3-Pstructure}
$P \setminus T \subseteq \agents \setminus S$; that is, every agent of $P$
that is not a recipient lies outside $S$.
\end{lemma}
\begin{proof}
Let $P_0 := T$ and let $P_{u+1}$ be $P_u$ together with all $i \notin S$
having an arc of $E$ into $P_u$; then $P = \bigcup_u P_u$. By induction on
$u$, $P_u \setminus T \subseteq \agents \setminus S$: this holds for
$P_0 \setminus T = \emptyset$, and every agent adjoined at any later step
lies outside $S$ by construction.
\end{proof}

\begin{lemma}
\label{lem:g3-budget}
$\subsidy \in \set{0,1}^n$ and $\sum_{i} \subsidy_i = \abs P \le n-1$.
\end{lemma}
\begin{proof}
That $\subsidy \in \set{0,1}^n$ and $\sum_i \subsidy_i = \abs P$ is immediate
from \eqref{eq:g3-subsidy}. By Lemma~\ref{lem:g3-Pstructure},
$P \subseteq T \cup (\agents \setminus S)$, and this union is disjoint
because $T \subseteq S$, so
$\abs P \le \abs T + \abs{\agents \setminus S} = r + n - \abs S$. By
Lemma~\ref{lem:g3-terminal}(a), $r \le \abs S - 1$, hence
$\abs P \le n - 1$.
\end{proof}

\begin{lemma}
\label{lem:g3-expensive}
Let $i \in S$, let $t \in T$ and let $e \in R$. Then
\[
  \cost_i(X_t \cup \set e) \ \ge\ \cost_i(X_i) + 1 .
\]
\end{lemma}
\begin{proof}
Note $t \in T \subseteq S$. By \eqref{eq:g3-pef},
$\cost_i(X_i) \le \cost_i(X_t)$.

If $\cost_i(X_i) < \cost_i(X_t)$, then $\cost_i(X_t) \ge \cost_i(X_i) + 1$ by
\eqref{eq:g3-gap}, and $\cost_i(X_t \cup \set e) \ge \cost_i(X_t)$ by
Observation~\ref{obs:g3-basic}.

Otherwise $\cost_i(X_i) = \cost_i(X_t)$. If $i = t$, then
$\cost_i(X_t \cup \set e) = \cost_i(X_i) + \cost_i(e \mid X_i) =
\cost_i(X_i) + 1$ by Lemma~\ref{lem:g3-terminal}(b). If $i \ne t$, then
$(i,t) \in E$ by \eqref{eq:g3-graph} with $i,t \in S$, so
$\cost_i(e \mid X_t) = 1$ by Lemma~\ref{lem:g3-terminal}(c), giving
$\cost_i(X_t \cup \set e) = \cost_i(X_t) + 1 = \cost_i(X_i) + 1$.
\end{proof}

\subsection{Envy-freeness}
\label{sec:g3-ef}

\begin{theorem}
\label{thm:g3-completion}
Let $X$, $R$, $r \ge 1$ and $S$ be as in \Cref{sec:g3-terminal}, let $\alloc$
be as in \eqref{eq:g3-alloc} and let $\subsidy$ be as in
\eqref{eq:g3-subsidy}. Then $(\alloc,\subsidy)$ is an envy-free solution.
\end{theorem}
\begin{proof}
Fix $i,j \in \agents$; we verify
$\cost_i(A_i) - \subsidy_i \le \cost_i(A_j) - \subsidy_j$. Throughout we use
that $X_j \subseteq A_j$ for every $j$, so
\begin{equation}
\label{eq:g3-mono}
  \cost_i(A_j) \ \ge\ \cost_i(X_j) \ \ge\ \cost_i(X_i)
\end{equation}
by Observation~\ref{obs:g3-basic} and \eqref{eq:g3-pef}. Note also that
$T \subseteq P$, so $i \notin P$ implies $i \notin T$ and hence $A_i = X_i$;
and $j \notin P$ implies $A_j = X_j$ and $\subsidy_j = 0$.

\smallskip
\noindent\textbf{Case 1: $i \in T$.}
Then $\subsidy_i = 1$ and, by Lemma~\ref{lem:g3-terminal}(b),
$\cost_i(A_i) = \cost_i(X_i) + 1$; hence
$\cost_i(A_i) - \subsidy_i = \cost_i(X_i)$. Note $i \in T \subseteq S$.
\begin{itemize}
  \item[(1a)] $j \in T$. Then $\subsidy_j = 1$ and $A_j = X_j \cup \set{e}$
    for some $e \in R$, so $\cost_i(A_j) \ge \cost_i(X_i) + 1$ by
    Lemma~\ref{lem:g3-expensive}. Hence
    $\cost_i(A_j) - \subsidy_j \ge \cost_i(X_i)$.
  \item[(1b)] $j \in P \setminus T$. Then $\subsidy_j = 1$ and $A_j = X_j$.
    By Lemma~\ref{lem:g3-Pstructure}, $j \notin S$; since $i \in S$,
    Lemma~\ref{lem:g3-terminal}(d) gives $\cost_i(X_i) < \cost_i(X_j)$, so
    $\cost_i(X_j) \ge \cost_i(X_i) + 1$ by \eqref{eq:g3-gap}. Hence
    $\cost_i(A_j) - \subsidy_j = \cost_i(X_j) - 1 \ge \cost_i(X_i)$.
  \item[(1c)] $j \notin P$. Then $\subsidy_j = 0$ and $A_j = X_j$, so
    $\cost_i(A_j) - \subsidy_j = \cost_i(X_j) \ge \cost_i(X_i)$ by
    \eqref{eq:g3-pef}.
\end{itemize}
In every sub-case $\cost_i(A_j) - \subsidy_j \ge \cost_i(X_i) =
\cost_i(A_i) - \subsidy_i$.

\smallskip
\noindent\textbf{Case 2: $i \in P \setminus T$.}
Then $\subsidy_i = 1$ and $A_i = X_i$, so
$\cost_i(A_i) - \subsidy_i = \cost_i(X_i) - 1$. For every $j$,
$\subsidy_j \le 1$ and $\cost_i(A_j) \ge \cost_i(X_i)$ by
\eqref{eq:g3-mono}, hence
$\cost_i(A_j) - \subsidy_j \ge \cost_i(X_i) - 1 = \cost_i(A_i) - \subsidy_i$.

\smallskip
\noindent\textbf{Case 3: $i \notin P$.}
Then $\subsidy_i = 0$ and $A_i = X_i$, so
$\cost_i(A_i) - \subsidy_i = \cost_i(X_i)$.
\begin{itemize}
  \item[(3a)] $j \notin P$. Then $\subsidy_j = 0$ and $A_j = X_j$, so
    $\cost_i(A_j) - \subsidy_j = \cost_i(X_j) \ge \cost_i(X_i)$ by
    \eqref{eq:g3-pef}.
  \item[(3b)] $j \in T$. Then $\subsidy_j = 1$ and $A_j = X_j \cup \set e$
    for some $e \in R$. We claim $\cost_i(A_j) \ge \cost_i(X_i) + 1$.
    If $i \in S$ this is Lemma~\ref{lem:g3-expensive}. If $i \notin S$,
    suppose $\cost_i(X_i) = \cost_i(X_j)$; then $i \ne j$ (as $j \in T
    \subseteq P$ and $i \notin P$), so $(i,j) \in E$ by
    \eqref{eq:g3-graph}, and $j \in T \subseteq P$, so (P2) of
    Definition~\ref{def:g3-P} places $i \in P$, contradicting $i \notin P$.
    Hence $\cost_i(X_i) < \cost_i(X_j)$, so $\cost_i(X_j) \ge \cost_i(X_i)+1$
    by \eqref{eq:g3-gap}, and $\cost_i(A_j) \ge \cost_i(X_j)$ by
    Observation~\ref{obs:g3-basic}. In both cases
    $\cost_i(A_j) - \subsidy_j \ge \cost_i(X_i)$.
  \item[(3c)] $j \in P \setminus T$. Then $\subsidy_j = 1$ and $A_j = X_j$.
    We claim $\cost_i(X_j) \ge \cost_i(X_i) + 1$. If $i \notin S$, suppose
    $\cost_i(X_i) = \cost_i(X_j)$; then $i \ne j$ (as $j \in P$ and
    $i \notin P$), so $(i,j) \in E$, and (P2) places $i \in P$, a
    contradiction; hence $\cost_i(X_i) < \cost_i(X_j)$ and \eqref{eq:g3-gap}
    applies. If $i \in S$, then $j \notin S$ by
    Lemma~\ref{lem:g3-Pstructure}, and Lemma~\ref{lem:g3-terminal}(d) gives
    $\cost_i(X_i) < \cost_i(X_j)$, so again \eqref{eq:g3-gap} applies. Hence
    $\cost_i(A_j) - \subsidy_j = \cost_i(X_j) - 1 \ge \cost_i(X_i)$.
\end{itemize}
In every sub-case $\cost_i(A_j) - \subsidy_j \ge \cost_i(X_i) =
\cost_i(A_i) - \subsidy_i$.

\smallskip
The three cases are exhaustive, since $\agents$ is the disjoint union of $T$,
$P \setminus T$ and $\agents \setminus P$.
\end{proof}

\subsection{Conjecture 2 for every number of agents}
\label{sec:g3-main}

\begin{theorem}
\label{thm:g3-main}
Conjecture~\ref{conj:main} holds for every $n$: for all dichotomous cost
functions $\cost_1,\dots,\cost_n$ on a finite chore set $\items$ there exist
an allocation $\alloc$ and a subsidy $\subsidy \in \set{0,1}^n$ with
$\sum_i \subsidy_i \le n-1$ such that $(\alloc,\subsidy)$ is an envy-free
solution, and both can be computed in polynomial time.
\end{theorem}
\begin{proof}
Run the algorithm of Theorem~\ref{thm:g3-twyz} and let $X$, $R$, $r$, $S$ be
as in \Cref{sec:g3-terminal}.

If $r = 0$ then $X$ is an allocation satisfying \eqref{eq:g3-pef}, so
$(X, \mathbf 0)$ is an envy-free solution with $\sum_i \subsidy_i = 0 \le n-1$.

If $r \ge 1$, form $T$, $\alloc$ and $\subsidy$ as in
\Cref{sec:g3-construction}. Theorem~\ref{thm:g3-completion} shows
$(\alloc,\subsidy)$ is an envy-free solution and
Lemma~\ref{lem:g3-budget} gives $\subsidy \in \set{0,1}^n$ with
$\sum_i \subsidy_i \le n-1$.

For the running time: the algorithm of Theorem~\ref{thm:g3-twyz} is
polynomial and returns $X$, $R$ and $S$. Constructing $\alloc$ takes $r \le
n-1$ assignments. The graph $G$ has $n$ vertices and is determined by the
$n^2$ values $\cost_i(X_j)$, each one oracle call. Computing $P$ is a
backward reachability search in $G$ from $T$, restricted to vertices outside
$S$, which takes time $O(n^2)$. Hence the work after
Theorem~\ref{thm:g3-twyz} is polynomial, and so is the whole procedure.
\end{proof}

\begin{remark}
\label{rem:g3-dependency}
Theorem~\ref{thm:g3-main} uses Theorem~\ref{thm:g3-twyz} together with the
three properties of its halting state recorded in
Lemma~\ref{lem:g3-terminal}, which are read off the rules (R1)--(R3) of
\Cref{sec:g3-terminal}; the statement of Theorem~\ref{thm:g3-twyz} alone does
not supply them. Beyond \eqref{eq:g3-dich}, no assumption is made on the cost
functions: they need not be additive or submodular, and no relation between
different agents' cost functions is assumed. Envy-freeness is verified
directly from \eqref{eq:g3-ef} in Theorem~\ref{thm:g3-completion}, and
Lemma~\ref{lem:g3-terminal}(a) is used only in
Lemma~\ref{lem:g3-budget}; the conclusion of
Theorem~\ref{thm:g3-completion} holds whenever $r \le \abs S$.
\end{remark}
